\clearpage
\VerdierPhysicalPageAnchor{0234}%
\begingroup\small\setlength{\parskip}{0pt}\setlength{\emergencystretch}{1.5em}
\setlength{\abovedisplayskip}{4pt plus 2pt minus 2pt}%
\setlength{\belowdisplayskip}{4pt plus 2pt minus 2pt}%

\VerdierSemanticAnchor{ch4d-4-1-9-cont}
\noindent 以
$\bigl({}^{1}\!Y(p,q),{}^{1}\!\delta(p,q,r)\bigr)$（相应地，
$\bigl({}^{2}\!Y(p,q),{}^{2}\!\delta(p,q,r)\bigr)$）表示由复形 $X$ 按第一次数
（相应地，第二次数）滤过而定义的、取值于 $K(\mathcal A)$ 的谱对象
(4.1.7)。由定义立即得到下列关系：

\VerdierSemanticAnchor{ch4d-4-1-9-1}
\begin{equation}
\left\{
\begin{aligned}
{}^{1}\!Y(q-1,q)&=X^{-q,\bullet}[q],\\
{}^{2}\!Y(q-1,q)&=X^{\bullet,-q}[q],\\
{}^{1}\!\delta(q-1,q,q+1)&=-\bigl(d_1^{-q-1,\bullet}\bigr)[q+1],\\
{}^{2}\!\delta(q-1,q,q+1)&=-\bigl(d_2^{\bullet,-q-1}\bigr)[q+1].
\end{aligned}
\right.
\tag{4.1.9.1}
\end{equation}
这些公式将在 (4.6) 中使用。

\VerdierSemanticAnchor{ch4d-4-2}
\medskip\noindent\textbf{4.2. 与 $\operatorname{comp}(\mathcal A)$
的滤过对象相联系的、取值于 $D(\mathcal A)$ 的谱对象。}

\VerdierSemanticAnchor{ch4d-4-2-scope}
\begin{center}\itshape
在本目以及本节以下各目中，$\mathcal A$ 表示阿贝尔范畴。
\end{center}

\VerdierSemanticAnchor{ch4d-4-2-1}
\noindent\textbf{4.2.1.}\quad
设 $X$ 是 $\operatorname{comp}(\mathcal A)$ 的对象，并设
$X(q)$，$q\in\mathbb Z$，是复形 $X$ 的
$\widetilde{\mathbb Z}$ 型递增滤过：

\VerdierSemanticAnchor{ch4d-4-2-1-1}
\begin{equation}
0=X(-\infty)\subset\cdots\subset X(q)\subset X(q+1)
\subset\cdots\subset X(+\infty)=X.
\tag{4.2.1.1}
\end{equation}

\VerdierSemanticAnchor{ch4d-4-2-1-associated-object}
由 (4.1.1) 所回顾的构造，得到一个取值于 $K(\mathcal A)$ 的谱对象
$\bigl(\mathcal X(p,q),\delta(p,q,r)\bigr)$；再应用典范函子
$K(\mathcal A)\longrightarrow D(\mathcal A)$，便得到一个取值于
$D(\mathcal A)$ 的谱对象，仍滥用记号写作
$\bigl(\mathcal X(p,q),\delta(p,q,r)\bigr)$。下面将描述一个取值于
$D(\mathcal A)$ 且与这个谱对象同构的谱对象。

\VerdierSemanticAnchor{ch4d-4-2-2}
\noindent\textbf{4.2.2.}\quad
因此，设
$0=X(-\infty)\subset\cdots\subset X(q)\subset\cdots\subset
X(+\infty)=X$ 是 $\operatorname{comp}(\mathcal A)$ 的滤过对象。
令：

\VerdierSemanticAnchor{ch4d-4-2-2-1}
\begin{equation}
X(p,q)=X(q)/X(p),
\tag{4.2.2.1}
\end{equation}
其中 $(p,q)$ 是 $\mathcal{F}\ell_1(\widetilde{\mathbb Z})$ 的任意元素。
（这里的商显然是在阿贝尔范畴 $\operatorname{comp}(\mathcal A)$ 的意义下
取得的。）记

\VerdierSemanticAnchor{ch4d-4-2-2-2}
\begin{equation}
\varphi\bigl((p,q),(p',q')\bigr)\colon
X(p,q)\longrightarrow X(p',q'),
\tag{4.2.2.2}
\end{equation}
\endgroup
